\begin{pastexam}[label={px:nccu-109-dsa-03}]{政大資科所 109 學年度 第 3 題}
（程式設計）

(20 points) Given two binary trees, design an algorithm for checking if they are identical in terms
of both the structure and the values. The node in the binary tree is defined as follows. Implement
the ``equal'' function in C or pseudo code.
\begin{minted}[fontsize=\small]{c}
struct node {
    int value;
    struct node *left, *right;
};
bool equal(struct node *tree1, struct node *tree2)
{
    ....
}
\end{minted}
\end{pastexam}

\begin{pastsoln}
\textbf{想法}。 兩棵二元樹相同，若且唯若：兩棵都是空的；或者兩棵都不空、根的值相同、左子樹相同、右子樹
相同。 這正是二元樹的遞迴定義（資結書第~7.3~節的走訪也是這樣遞迴），直接寫成遞迴函式：
\begin{code}{c}{判斷兩棵二元樹是否相同}
bool equal(struct node *tree1, struct node *tree2)
{
    if (tree1 == NULL || tree2 == NULL)
        return tree1 == tree2;          /* 兩個都空才相同 */
    return tree1->value == tree2->value
        && equal(tree1->left, tree2->left)
        && equal(tree1->right, tree2->right);
}
\end{code}
（\texttt{bool} 要 \texttt{\#include <stdbool.h>}。）

\textbf{正確性}。 對兩棵樹的節點數總和 $s$ 做歸納。
\begin{itemize}
\item 至少一棵是空的：兩棵都空時相同，回傳 \texttt{true}；只有一棵空時結構不同，回傳 \texttt{false}。
第 3--4 行用 \texttt{tree1 == tree2} 一次處理這兩種情形（兩個都是 \texttt{NULL} 才相等）。
\item 兩棵都不空：它們相同，若且唯若根的值相同、而且左子樹和左子樹相同、右子樹和右子樹相同。 子樹的節點數
總和比 $s$ 小，由歸納假設遞迴呼叫回傳正確的結果，第 5--7 行把三個條件用 \texttt{\&\&} 接起來。
\end{itemize}
\texttt{\&\&} 是短路求值：根的值不同就不會再遞迴，左子樹不同就不會再比右子樹，只會更快，不影響結果。

\textbf{複雜度}。 設兩棵樹分別有 $n_1$、$n_2$ 個節點。 每次呼叫做常數的工作；只有兩邊都不空時才會往下遞迴，
所以被呼叫的節點對最多 $\min(n_1, n_2)$ 個，加上碰到 \texttt{NULL} 的呼叫（每個節點對最多產生兩個），
總共 $O(\min(n_1, n_2))$。 最壞（兩棵相同）是 $\Theta(n)$。 遞迴深度不超過較矮那棵的高度加一，最壞
$O(n)$ 的堆疊空間。

\textbf{常見的錯}。 只比較兩棵樹的中序（或前序）走訪序列不夠：只有 $1$、$2$ 兩個值時，「根 $1$、右孩子 $2$」
和「根 $2$、左孩子 $1$」的中序都是 $1, 2$，結構卻不同。 一定要連空指標的位置一起比，上面的函式就是在做這件事。
\end{pastsoln}
